\section{Sous-préschémas et morphismes d'immersion}
\label{section:I.4-fr}

\subsection{Sous-préschémas.}
\label{subsection:I.4.1-fr}

\begin{env}[4.1.1]
\label{I.4.1.1-fr}
Comme la notion de faisceau quasi-cohérent
(\hyperref[0.5.1.3-fr]{0, 5.1.3}) est locale, un
$\mathscr{O}_X$-Module quasi-cohérent $\mathscr{F}$ sur un préschéma $X$ peut
être défini par la condition que, pour tout ouvert affine $V$ de $X$,
$\mathscr{F}|V$ est isomorphe à un faisceau associé à un
$\Gamma(V,\mathscr{O}_X)$-module (\hyperref[I.1.4.1-fr]{1.4.1}). Il est clair
que, sur un préschéma $X$, le faisceau structural $\mathscr{O}_X$ est
quasi-cohérent et que noyaux, conoyaux, images d'homomorphismes de
$\mathscr{O}_X$-Modules quasi-cohérents, limites inductives et sommes directes
de $\mathscr{O}_X$-Modules quasi-cohérents sont quasi-cohérents
(\hyperref[I.1.3.7-fr]{1.3.7} et \hyperref[I.1.3.9-fr]{1.3.9}).
\end{env}

\begin{proposition}[4.1.2]
\label{I.4.1.2-fr}
Soient $X$ un préschéma, $\mathscr{I}$ un faisceau quasi-cohérent d'idéaux
dans $\mathscr{O}_X$. Le support $Y$ du faisceau
$\mathscr{O}_X/\mathscr{I}$ est alors fermé, et si on désigne par
$\mathscr{O}_Y$ la restriction de $\mathscr{O}_X/\mathscr{I}$ à $Y$,
$(Y,\mathscr{O}_Y)$ est un préschéma.
\end{proposition}

\oldpage[I]{120}
Il suffit évidemment (\hyperref[I.2.1.3-fr]{2.1.3}) de considérer le cas où
$X$ est un schéma affine, et de montrer que dans ce cas $Y$ est fermé dans
$X$ et est un \emph{schéma affine}. En effet, si
$X=\operatorname{Spec}(A)$, on a $\mathscr{O}_X=\widetilde{A}$ et
$\mathscr{I}=\widetilde{\mathfrak{I}}$, où $\mathfrak{I}$ est un idéal de
$A$ (\hyperref[I.1.4.1-fr]{1.4.1})~; $Y$ est alors égal à la partie fermée
$V(\mathfrak{I})$ de $X$ et s'identifie au spectre premier de l'anneau
$B=A/\mathfrak{I}$ (\hyperref[I.1.1.11-fr]{1.1.11})~; de plus, si $\varphi$
est l'homomorphisme canonique $A\to B=A/\mathfrak{I}$, l'image directe
${}^a\varphi_*(\widetilde{B})$ s'identifie canoniquement au faisceau
$\widetilde{A}/\widetilde{\mathfrak{I}}=\mathscr{O}_X/\mathscr{I}$
(\hyperref[I.1.6.3-fr]{1.6.3} et \hyperref[I.1.3.9-fr]{1.3.9}), ce qui achève
la démonstration.

Nous dirons que $(Y,\mathscr{O}_Y)$ est le \emph{sous-préschéma} de
$(X,\mathscr{O}_X)$ \emph{défini par le faisceau d'idéaux $\mathscr{I}$}~;
c'est un cas particulier de la notion générale de \emph{sous-préschéma}~:

\begin{definition}[4.1.3]
\label{I.4.1.3-fr}
On dit qu'un espace annelé $(Y,\mathscr{O}_Y)$ est un
\emph{sous-préschéma} d'un préschéma $(X,\mathscr{O}_X)$ si~:
\begin{enumerate}
  \item[$1^\circ$] $Y$ est un sous-espace localement fermé de $X$~;
  \item[$2^\circ$] Si $U$ désigne le plus grand ouvert de $X$ contenant $Y$
  et tel que $Y$ soit fermé dans $U$ (autrement dit, le complémentaire dans
  $X$ de la frontière de $Y$ par rapport à $\overline{Y}$),
  $(Y,\mathscr{O}_Y)$ est un sous-préschéma de
  $(U,\mathscr{O}_X|U)$ défini par un faisceau quasi-cohérent d'idéaux de
  $\mathscr{O}_X|U$.
\end{enumerate}
On dit que le sous-préschéma $(Y,\mathscr{O}_Y)$ de $(X,\mathscr{O}_X)$ est
\emph{fermé} si $Y$ est fermé dans $X$ (auquel cas $U=X$).
\end{definition}

Il résulte aussitôt de cette définition et de
(\hyperref[I.4.1.2-fr]{4.1.2}) que les sous-préschémas fermés de $X$ sont en
\emph{correspondance biunivoque canonique} avec les faisceaux d'idéaux
quasi-cohérents $\mathscr{J}$ de $\mathscr{O}_X$, car le fait que deux tels
faisceaux $\mathscr{J}$, $\mathscr{J}'$ aient même support (fermé) $Y$ et
soient tels que les restrictions de $\mathscr{O}_X/\mathscr{J}$ et de
$\mathscr{O}_X/\mathscr{J}'$ à $Y$ soient identiques, entraîne aussitôt que
$\mathscr{J}'=\mathscr{J}$.

\begin{env}[4.1.4]
\label{I.4.1.4-fr}
Soient $(Y,\mathscr{O}_Y)$ un sous-préschéma de $X$, $U$ le plus grand ouvert
de $X$ contenant $Y$ et dans lequel $Y$ soit fermé, $V$ un ouvert de $X$
contenu dans $U$~; alors $V\cap Y$ est fermé dans $V$. En outre, si $Y$ est
défini par le faisceau quasi-cohérent $\mathscr{J}$ d'idéaux de
$\mathscr{O}_X|U$, $\mathscr{J}|V$ est un faisceau quasi-cohérent d'idéaux de
$\mathscr{O}_X|V$, et il est immédiat que le préschéma induit par $Y$ sur
$Y\cap V$ est le sous-préschéma fermé de $V$ défini par le faisceau d'idéaux
$\mathscr{J}|V$. Inversement~:
\end{env}

\begin{proposition}[4.1.5]
\label{I.4.1.5-fr}
Soit $(Y,\mathscr{O}_Y)$ un espace annelé tel que $Y$ soit un sous-espace de
$X$ et qu'il existe un recouvrement $(V_\alpha)$ de $Y$ par des ouverts de
$X$ tels que pour tout $\alpha$, $Y\cap V_\alpha$ soit fermé dans $V_\alpha$
et que l'espace annelé
$(Y\cap V_\alpha,\mathscr{O}_Y|(Y\cap V_\alpha))$ soit un sous-préschéma
fermé du préschéma induit sur $V_\alpha$ par $X$. Alors
$(Y,\mathscr{O}_Y)$ est un sous-préschéma de $X$.
\end{proposition}

L'hypothèse implique que $Y$ est localement fermé dans $X$ et que le plus
grand ouvert $U$ contenant $Y$ et dans lequel $Y$ est fermé contient tous les
$V_\alpha$~; on peut donc se ramener au cas où $U=X$ et $Y$ est fermé dans
$X$. On définit alors un faisceau quasi-cohérent d'idéaux $\mathscr{J}$ de
$\mathscr{O}_X$ en prenant pour $\mathscr{J}|V_\alpha$ le faisceau d'idéaux de
$\mathscr{O}_X|V_\alpha$ qui définit le sous-préschéma fermé
$(Y\cap V_\alpha,\mathscr{O}_Y|(Y\cap V_\alpha))$ et pour tout ouvert $W$ de
$X$ ne rencontrant pas $Y$, $\mathscr{J}|W=\mathscr{O}_X|W$. On vérifie
immédiatement en vertu de (\hyperref[I.4.1.3-fr]{4.1.3}) et
(\hyperref[I.4.1.4-fr]{4.1.4}) qu'il existe un faisceau d'idéaux $\mathscr{J}$
et un seul satisfaisant à ces conditions et qu'il définit le sous-préschéma
fermé $(Y,\mathscr{O}_Y)$.

En particulier, le préschéma \emph{induit} par $X$ sur un \emph{ouvert} de
$X$ est un \emph{sous-préschéma} de $X$.

\begin{proposition}[4.1.6]
\label{I.4.1.6-fr}
Un sous-préschéma (resp. un sous-préschéma fermé) d'un
sous-\oldpage[I]{121}préschéma (resp. sous-préschéma fermé) de $X$
s'identifie canoniquement à un sous-préschéma (resp. sous-préschéma fermé) de
$X$.
\end{proposition}

Une partie localement fermée d'un sous-espace localement fermé de $X$ étant
un sous-espace localement fermé de $X$, il est clair
(\hyperref[I.4.1.5-fr]{4.1.5}) que la question est locale et qu'on peut donc
supposer $X$ affine~; la proposition résulte alors de l'identification
canonique de $A/\mathfrak{J}'$ et de
$(A/\mathfrak{J})/(\mathfrak{J}'/\mathfrak{J})$ lorsque
$\mathfrak{J}$, $\mathfrak{J}'$ sont deux idéaux d'un anneau $A$ tels que
$\mathfrak{J}\subset\mathfrak{J}'$.

Nous ferons toujours par la suite l'identification précédente.

\begin{env}[4.1.7]
\label{I.4.1.7-fr}
Soit $Y$ un sous-préschéma d'un préschéma $X$, et désignons par $\psi$
l'injection canonique $Y\to X$ des \emph{espaces sous-jacents}~; on sait que
l'image réciproque $\psi^*(\mathscr{O}_X)$ est la restriction
$\mathscr{O}_X|Y$ (\hyperref[0.3.7.1-fr]{0, 3.7.1}). Si, pour tout $y\in Y$,
on désigne par $\omega_y$ l'homomorphisme canonique
$(\mathscr{O}_X)_y\to(\mathscr{O}_Y)_y$, ces homomorphismes sont les
restrictions aux fibres d'un homomorphisme \emph{surjectif} $\omega$ de
faisceaux d'anneaux $\mathscr{O}_X|Y\to\mathscr{O}_Y$~: il suffit en effet de
le vérifier localement sur $Y$, c'est-à-dire que l'on peut supposer $X$
affine et le sous-préschéma $Y$ fermé~; si dans ce cas $\mathscr{I}$ est le
faisceau d'idéaux dans $\mathscr{O}_X$ qui définit $Y$, les $\omega_y$ ne sont
autres que les restrictions aux fibres de l'homomorphisme
$\mathscr{O}_X|Y\to(\mathscr{O}_X/\mathscr{I})|Y$.

Nous avons donc défini un \emph{monomorphisme d'espaces annelés}
(\hyperref[0.4.1.1-fr]{0, 4.1.1}) $j=(\psi,\omega^b)$ qui est évidemment un
morphisme $Y\to X$ de préschémas (\hyperref[I.2.2.1-fr]{2.2.1}), et que nous
appellerons le \emph{morphisme d'injection canonique}.

Si $f:X\to Z$ est un morphisme, nous dirons encore que le morphisme composé
$Y\xrightarrow{j}X\xrightarrow{f}Z$ est la \emph{restriction} de $f$ au
sous-préschéma $Y$.
\end{env}

\begin{env}[4.1.8]
\label{I.4.1.8-fr}
Conformément aux définitions générales (T, I, 1.1), nous dirons qu'un
morphisme de préschémas $f:Z\to X$ est \emph{majoré} par le morphisme
d'injection $j:Y\to X$ d'un sous-préschéma $Y$ de $X$ si $f$ se factorise en
$Z\xrightarrow{g}Y\xrightarrow{j}X$, où $g$ est un morphisme de préschémas~;
$g$ est nécessairement \emph{unique} puisque $j$ est un monomorphisme.
\end{env}

\begin{proposition}[4.1.9]
\label{I.4.1.9-fr}
Pour qu'un morphisme $f:Z\to X$ soit \emph{majoré par un morphisme
d'injection} $j:Y\to X$, il faut et il suffit que $f(Z)\subset Y$ et que,
pour tout $z\in Z$, si on pose $y=f(z)$, l'homomorphisme
$(\mathscr{O}_X)_y\to\mathscr{O}_z$ correspondant à $f$ se factorise en
$(\mathscr{O}_X)_y\to(\mathscr{O}_Y)_y\to\mathscr{O}_z$ (ou, ce qui revient
au même, que le noyau de $(\mathscr{O}_X)_y\to\mathscr{O}_z$ contienne celui
de $(\mathscr{O}_X)_y\to(\mathscr{O}_Y)_y$).
\end{proposition}

Les conditions sont évidemment nécessaires. Pour voir qu'elles sont
suffisantes, on peut se ramener au cas où $Y$ est un sous-préschéma
\emph{fermé} de $X$, en remplaçant au besoin $X$ par un ouvert $U$ tel que
$Y$ soit fermé dans $U$ (\hyperref[I.4.1.3-fr]{4.1.3})~; $Y$ est alors défini
par un faisceau quasi-cohérent $\mathscr{I}$ d'idéaux de $\mathscr{O}_X$.
Posons $f=(\psi,\theta)$, et soit $\mathscr{J}$ le faisceau d'idéaux de
$\psi^*(\mathscr{O}_X)$, noyau de
$\theta^\sharp:\psi^*(\mathscr{O}_X)\to\mathscr{O}_Z$~; compte tenu des
propriétés du foncteur $\psi^*$ (\hyperref[0.3.7.2-fr]{0, 3.7.2}),
l'hypothèse entraîne que pour tout $z\in Z$, on a
$(\psi^*(\mathscr{I}))_z\subset\mathscr{J}_z$, et par suite
$\psi^*(\mathscr{I})\subset\mathscr{J}$. Par suite, $\theta^\sharp$ se
factorise en
\[
  \psi^*(\mathscr{O}_X)\longrightarrow
  \psi^*(\mathscr{O}_X)/\psi^*(\mathscr{I})
  =\psi^*(\mathscr{O}_X/\mathscr{I})
  \xrightarrow{\omega}\mathscr{O}_Z,
\]
la première flèche étant l'homomorphisme canonique. Soit $\psi'$ l'application
continue $Z\to Y$ coïncidant avec $\psi$~; il est clair que l'on a
${\psi'}^*(\mathscr{O}_Y)=\psi^*(\mathscr{O}_X/\mathscr{I})$~; d'autre part,
$\omega$ est évidemment un homomorphisme local, donc
$g'=(\psi',\omega^b)$ est un morphisme $Z\to Y$
\oldpage[I]{122}
de préschémas (\hyperref[I.2.2.1-fr]{2.2.1}), qui, en vertu de ce qui précède,
est tel que $f=j\circ g$, d'où la proposition.

\begin{corollary}[4.1.10]
\label{I.4.1.10-fr}
Pour qu'un morphisme d'injection $Z\to X$ soit \emph{majoré par le morphisme
d'injection} $Y\to X$, il faut et il suffit que $Z$ soit un sous-préschéma de
$Y$.
\end{corollary}

On écrit alors $Z\leq Y$, et cette relation est évidemment une
\emph{relation d'ordre} dans l'ensemble des sous-préschémas de $X$.

\subsection{Morphismes d'immersion}
\label{subsection:I.4.2-fr}

\begin{definition}[4.2.1]
\label{I.4.2.1-fr}
On dit qu'un morphisme $f:Y\to X$ est une \emph{immersion} (resp. une
\emph{immersion fermée}, une \emph{immersion ouverte}) s'il se factorise en
$Y\xrightarrow{g}Z\xrightarrow{j}X$, où $g$ est un isomorphisme, $Z$ un
sous-préschéma de $X$ (resp. un sous-préschéma fermé, un sous-préschéma induit
sur un ouvert) et $j$ le morphisme d'injection.
\end{definition}

Le sous-préschéma $Z$ et l'isomorphisme $g$ sont alors déterminés de façon
\emph{unique}, car si $Z'$ est un second sous-préschéma de $X$, $j'$ l'injection
$Z'\to X$ et $g'$ un isomorphisme $Y\to Z'$ tel que
$j\circ g=j'\circ g'$, on en déduit $j'=j\circ g\circ {g'}^{-1}$, d'où
$Z'\leq Z$ (\hyperref[I.4.1.10-fr]{4.1.10}), et on montre de même que
$Z\leq Z'$, donc $Z'=Z$, et comme $j$ est un monomorphisme de préschémas,
$g'=g$.

On dit que $f=j\circ g$ est la \emph{factorisation canonique} de l'immersion
$f$, et le sous-préschéma $Z$ et l'isomorphisme $g$ sont dits
\emph{associés à $f$}.

Il est clair qu'une immersion est un \emph{monomorphisme} de préschémas
(\hyperref[I.4.1.7-fr]{4.1.7}) et \emph{a fortiori} un morphisme
\emph{radiciel} (\hyperref[I.3.5.4-fr]{3.5.4}).

\begin{proposition}[4.2.2]
\label{I.4.2.2-fr}
\begin{enumerate}
  \item[a)] Pour qu'un morphisme $f=(\psi,\theta):Y\to X$ soit une
  \emph{immersion ouverte}, il faut et il suffit que $\psi$ soit un
  homéomorphisme de $Y$ sur une partie ouverte de $X$, et que pour tout
  $y\in Y$, l'homomorphisme
  $\theta_y^\sharp:\mathscr{O}_{\psi(y)}\to\mathscr{O}_y$ soit bijectif.
  \item[b)] Pour qu'un morphisme $f=(\psi,\theta):Y\to X$ soit une
  \emph{immersion} (resp. une \emph{immersion fermée}), il faut et il suffit
  que $\psi$ soit un homéomorphisme de $Y$ sur une partie localement fermée
  (resp. fermée) de $X$, et que pour tout $y\in Y$, l'homomorphisme
  $\theta_y^\sharp:\mathscr{O}_{\psi(y)}\to\mathscr{O}_y$ soit surjectif.
\end{enumerate}
\end{proposition}

\begin{enumerate}
  \item[a)] Les conditions sont évidemment nécessaires. Inversement, si elles
  sont remplies, il est clair que $\theta^\sharp$ est un isomorphisme de
  $\mathscr{O}_Y$ sur $\psi^*(\mathscr{O}_X)$, et
  $\psi^*(\mathscr{O}_X)$ est le faisceau déduit par transport de structure au
  moyen de $\psi^{-1}$ à partir de $\mathscr{O}_X|\psi(Y)$~; d'où la
  conclusion.
  \item[b)] Les conditions étant évidemment nécessaires, prouvons qu'elles
  sont suffisantes. Considérons d'abord le cas particulier où l'on suppose que
  $X$ est un schéma affine et que $Z=\psi(Y)$ est \emph{fermé} dans $X$. On
  sait (\hyperref[0.3.4.6-fr]{0, 3.4.6}) que $\psi_*(\mathscr{O}_Y)$ a alors
  pour support $Z$ et que si l'on désigne par $\mathscr{O}'_Z$ sa restriction
  à $Z$, l'espace annelé $(Z,\mathscr{O}'_Z)$ se déduit de
  $(Y,\mathscr{O}_Y)$ par transport de structure au moyen de
  l'homéomorphisme $\psi$, considéré comme application de $Y$ sur $Z$.
  Montrons que $f_*(\mathscr{O}_Y)=\psi_*(\mathscr{O}_Y)$ est un
  $\mathscr{O}_X$-Module \emph{quasi-cohérent}. En effet, pour tout
  $x\notin Z$, $\psi_*(\mathscr{O}_Y)$, restreint à un voisinage convenable de
  $x$, est nul. Si au contraire $x\in Z$, on a $x=\psi(y)$ pour un $y\in Y$
  bien déterminé~; soit $V$ un voisinage ouvert affine de $y$ dans $Y$~;
  $\psi(V)$ est alors ouvert dans $Z$, donc trace sur $Z$ d'un ouvert $U$ de
  $X$, et la restriction à $U$ de $\psi_*(\mathscr{O}_Y)$ est identique à la
  restriction à $U$ de l'image
\oldpage[I]{123}
  directe $(\psi_V)_*(\mathscr{O}_Y|V)$, où $\psi_V$ est la restriction de
  $\psi$ à $V$. Or, la restriction à $(V,\mathscr{O}_Y|V)$ du morphisme
  $(\psi,\theta)$ est un morphisme de ce préschéma dans
  $(X,\mathscr{O}_X)$, et par suite est de la forme
  $({}^a\varphi,\widetilde{\varphi})$, où $\varphi$ est un homomorphisme de
  l'anneau $A=\Gamma(X,\mathscr{O}_X)$ dans l'anneau
  $\Gamma(V,\mathscr{O}_Y)$ (\hyperref[I.1.7.3-fr]{1.7.3})~; on en conclut
  que $(\psi_V)_*(\mathscr{O}_Y|V)$ est un $\mathscr{O}_X$-Module
  quasi-cohérent (\hyperref[I.1.6.3-fr]{1.6.3}), ce qui prouve notre
  assertion, en raison du caractère local des faisceaux quasi-cohérents. En
  outre, l'hypothèse que $\psi$ est un homéomorphisme entraîne
  (\hyperref[0.3.4.5-fr]{0, 3.4.5}) que pour tout $y\in Y$, $\psi_y$ est un
  isomorphisme
  $(\psi_*(\mathscr{O}_Y))_{\psi(y)}\to\mathscr{O}_y$~; comme le diagramme
  \phantomsection\label{I.4.2.2.theta-diagram-fr}
  \[
    \xymatrix{
      \mathscr{O}_{\psi(y)}\ar[r]^{\theta_{\psi(y)}}
        \ar[d]_{\psi_y\circ\rho_{\psi(y)}} &
      (\psi_*(\mathscr{O}_Y))_{\psi(y)}\ar[d]^{\psi_y}\\
      (\psi^*(\mathscr{O}_X))_y\ar[r]^{\theta_y^\sharp} &
      \mathscr{O}_y
    }
  \]
  est commutatif et que les deux flèches verticales sont des isomorphismes
  (\hyperref[0.3.7.2-fr]{0, 3.7.2}), l'hypothèse que $\theta_y^\sharp$ est
  surjectif entraîne qu'il en est de même de $\theta_{\psi(y)}$. Comme le
  support de $\psi_*(\mathscr{O}_Y)$ est $Z=\psi(Y)$, $\theta$ est un
  homomorphisme \emph{surjectif} de $\mathscr{O}_X=\widetilde{A}$ dans le
  $\mathscr{O}_X$-Module quasi cohérent $f_*(\mathscr{O}_Y)$. Par suite, il
  existe un isomorphisme unique $\omega$ d'un faisceau quotient
  $\widetilde{A}/\widetilde{\mathfrak{J}}$ ($\mathfrak{J}$ idéal de $A$) sur
  $f_*(\mathscr{O}_Y)$ qui, composé avec l'homomorphisme canonique
  $\widetilde{A}\to\widetilde{A}/\widetilde{\mathfrak{J}}$, donne $\theta$
  (\hyperref[I.1.3.8-fr]{1.3.8})~; si $\mathscr{O}_Z$ désigne la restriction
  de $\widetilde{A}/\widetilde{\mathfrak{J}}$ à $Z$,
  $(Z,\mathscr{O}_Z)$ est un sous-préschéma de $(X,\mathscr{O}_X)$, et $f$ se
  factorise en l'injection canonique de ce sous-préschéma dans $X$, et
  l'isomorphisme $(\psi_0,\omega_0)$, où $\psi_0$ est $\psi$ considéré comme
  application de $Y$ sur $Z$, et $\omega_0$ la restriction de $\omega$ à
  $\mathscr{O}_Z$.

  Passons au cas général. Soit $U$ un ensemble ouvert affine dans $X$ tel que
  $U\cap\psi(Y)$ soit fermé dans $U$ et non vide. En restreignant $f$ au
  préschéma induit par $Y$ sur l'ouvert $\psi^{-1}(U)$, et en le considérant
  comme un morphisme de ce préschéma dans le préschéma induit par $X$ sur
  $U$, on est ramené au premier cas~; la restriction de $f$ à
  $\psi^{-1}(U)$ est donc une immersion fermée $\psi^{-1}(U)\to U$, se
  factorisant canoniquement en $j_U\circ g_U$, où $g_U$ est un isomorphisme
  du préschéma $\psi^{-1}(U)$ sur un sous-préschéma $Z_U$ de $U$, et $j_U$
  l'injection canonique $Z_U\to U$. Soit $V$ un second ouvert affine de $X$
  tel que $V\subset U$~; comme la restriction $Z'_V$ de $Z_U$ à $V$ est un
  sous-préschéma du préschéma $V$, la restriction de $f$ à $\psi^{-1}(V)$ se
  factorise en $j'_V\circ g'_V$, où $j'_V$ est l'injection canonique
  $Z'_V\to V$ et $g'_V$ un isomorphisme de $\psi^{-1}(V)$ sur $Z'_V$. Par
  l'unicité de la factorisation canonique d'une immersion
  (\hyperref[I.4.2.1-fr]{4.2.1}), on a nécessairement $Z'_V=Z_V$ et
  $g'_V=g_V$. On en conclut (\hyperref[I.4.1.5-fr]{4.1.5}) qu'il y a un
  sous-préschéma $Z$ de $X$ dont l'espace sous-jacent est $\psi(Y)$ et dont
  la restriction à chaque $U\cap\psi(Y)$ est $Z_U$~; les $g_U$ sont alors les
  restrictions aux $\psi^{-1}(U)$ d'un isomorphisme $g:Y\to Z$ tel que
  $f=j\circ g$, où $j$ est l'injection canonique $Z\to X$.
\end{enumerate}

\begin{corollary}[4.2.3]
\label{I.4.2.3-fr}
Soit $X$ un schéma affine. Pour qu'un morphisme
$f=(\psi,\theta):Y\to X$ soit une immersion fermée, il faut et il suffit que
$Y$ soit un schéma affine et que l'homomorphisme
$\Gamma(\psi):\Gamma(\mathscr{O}_X)\to\Gamma(\mathscr{O}_Y)$ soit surjectif.
\end{corollary}

\begin{corollary}[4.2.4]
\label{I.4.2.4-fr}
\begin{enumerate}
  \item[a)] Soient $f$ un morphisme $Y\to X$, $(V_\lambda)$ un recouvrement
  de $f(Y)$ par des ouverts de $X$. Pour que $f$ soit une immersion (resp. une
  immersion ouverte), il faut et il suffit
\oldpage[I]{124}
  que sa restriction à chacun des préschémas induits $f^{-1}(V_\lambda)$ soit
  une immersion (resp. une immersion ouverte) dans $V_\lambda$.
  \item[b)] Soient $f$ un morphisme $Y\to X$, $(V_\lambda)$ un recouvrement
  ouvert de $X$. Pour que $f$ soit une immersion fermée, il faut et il suffit
  que sa restriction à chacun des préschémas induits $f^{-1}(V_\lambda)$ soit
  une immersion fermée dans $V_\lambda$.
\end{enumerate}
\end{corollary}

Soit $f=(\psi,\theta)$~; dans le cas a), $\theta_y^\sharp$ est surjectif
(resp. bijectif) pour tout $y\in Y$, et dans le cas b) $\theta_y^\sharp$ est
surjectif pour tout $y\in Y$~; il suffit donc de vérifier que $\psi$, dans le
cas a), est un homéomorphisme de $Y$ sur une partie localement fermée (resp.
ouverte) de $X$, et dans le cas b), un homéomorphisme de $Y$ sur une partie
fermée de $X$. Or, $\psi$ est évidemment injectif et transforme tout voisinage
de $y$ dans $Y$ en un voisinage de $\psi(y)$ dans $\psi(Y)$ pour tout $y\in Y$,
en vertu de l'hypothèse~; dans le cas a), $\psi(Y)\cap V_\lambda$ est localement
fermé (resp. ouvert) dans $V_\lambda$, donc $\psi(Y)$ est localement fermé
(resp. ouvert) dans la réunion des $V_\lambda$, et \emph{a fortiori} dans $X$~;
dans le cas b), $\psi(Y)\cap V_\lambda$ est fermé dans $V_\lambda$, donc
$\psi(Y)$ est fermé dans $X$ puisque $X=\bigcup_\lambda V_\lambda$.

\begin{proposition}[4.2.5]
\label{I.4.2.5-fr}
Le composé de deux immersions (resp. de deux immersions ouvertes, de deux
immersions fermées) est une immersion (resp. une immersion ouverte, une
immersion fermée).
\end{proposition}

Cela résulte trivialement de (\hyperref[I.4.1.6-fr]{4.1.6}).

\subsection{Produit d'immersions.}
\label{subsection:I.4.3-fr}

\begin{proposition}[4.3.1]
\label{I.4.3.1-fr}
Soient $\alpha:X'\to X$, $\beta:Y'\to Y$ deux $S$-morphismes~; si $\alpha$ et
$\beta$ sont des immersions (resp. des immersions ouvertes, des immersions
fermées), $\alpha\times_S\beta$ est une immersion (resp. une immersion ouverte,
une immersion fermée). En outre, si $\alpha$ (resp. $\beta$) identifie $X'$
(resp. $Y'$) à un sous-préschéma $X''$ (resp. $Y''$) de $X$ (resp. $Y$),
$\alpha\times_S\beta$ identifie l'espace sous-jacent à $X'\times_S Y'$ au
sous-espace $p^{-1}(X'')\cap q^{-1}(Y'')$ de l'espace sous-jacent à
$X\times_S Y$, en désignant par $p$ et $q$ les projections de $X\times_S Y$
dans $X$ et $Y$ respectivement.
\end{proposition}

Compte tenu de la déf. (\hyperref[I.4.2.1-fr]{4.2.1}), on peut se restreindre
au cas où $X'$ et $Y'$ sont des sous-préschémas, $\alpha$ et $\beta$ les
morphismes d'injection. La proposition a déjà été établie pour les
sous-préschémas induits sur les ouverts (\hyperref[I.3.2.7-fr]{3.2.7})~; comme
tout sous-préschéma est un sous-préschéma fermé d'un préschéma induit sur un
ouvert (\hyperref[I.4.1.3-fr]{4.1.3}), on est ramené au cas où $X'$ et $Y'$ sont
des sous-préschémas \emph{fermés}.

Montrons d'abord qu'on peut supposer que $S$ est \emph{affine}. En effet, soit
$(S_\lambda)$ un recouvrement de $S$ formé d'ouverts affines~; si $\varphi$ et
$\psi$ sont les morphismes structuraux de $X$ et $Y$, soit
$X_\lambda=\varphi^{-1}(S_\lambda)$, $Y_\lambda=\psi^{-1}(S_\lambda)$. La
restriction $X'_\lambda$ (resp. $Y'_\lambda$) de $X'$ (resp. $Y'$) à
$X_\lambda\cap X'$ (resp. $Y_\lambda\cap Y'$) est un sous-préschéma fermé de
$X_\lambda$ (resp. $Y_\lambda$), les préschémas $X_\lambda$, $Y_\lambda$,
$X'_\lambda$, $Y'_\lambda$ peuvent être considérés comme des
$S_\lambda$-préschémas et les produits $X_\lambda\times_S Y_\lambda$ et
$X_\lambda\times_{S_\lambda}Y_\lambda$ (resp.
$X'_\lambda\times_S Y'_\lambda$ et
$X'_\lambda\times_{S_\lambda}Y'_\lambda$) sont identiques
(\hyperref[I.3.2.5-fr]{3.2.5}). Si la proposition est vraie lorsque $S$ est
affine, la restriction de $\alpha\times_S\beta$ à chacun des
$X'_\lambda\times_S Y'_\lambda$ sera donc une immersion
(\hyperref[I.3.2.7-fr]{3.2.7}). Comme le produit
$X'_\lambda\times_S Y'_\mu$ (resp. $X_\lambda\times_S Y_\mu$) s'identifie à
$(X'_\lambda\cap X'_\mu)\times_S(Y'_\lambda\cap Y'_\mu)$ (resp.
$(X_\lambda\cap X_\mu)\times_S(Y_\lambda\cap Y_\mu)$)
(\hyperref[I.3.2.6.4-fr]{3.2.6.4}), la restriction de
$\alpha\times_S\beta$
\oldpage[I]{125}
à chacun des $X'_\lambda\times_S Y'_\mu$ est encore une immersion~; il en est
par suite de même de $\alpha\times_S\beta$ en vertu de
(\hyperref[I.4.2.4-fr]{4.2.4}).

En second lieu, prouvons qu'on peut aussi supposer que $X$ et $Y$ sont
\emph{affines}. En effet, soit $(U_i)$ (resp. $(V_j)$) un recouvrement de $X$
(resp. $Y$) par des ouverts affines, et soit $X'_i$ (resp. $Y'_j$) la
restriction de $X'$ (resp. $Y'$) à $X'\cap U_i$ (resp. $Y'\cap V_j$), qui est
un sous-préschéma fermé de $U_i$ (resp. $V_j$)~; $U_i\times_S V_j$ s'identifie
à la restriction de $X\times_S Y$ à
$p^{-1}(U_i)\cap q^{-1}(V_j)$ (\hyperref[I.3.2.7-fr]{3.2.7})~; et de même, si
$p'$, $q'$ sont les projections de $X'\times_S Y'$,
$X'_i\times_S Y'_j$ s'identifie à la restriction de $X'\times_S Y'$ à
${p'}^{-1}(X'_i)\cap{q'}^{-1}(Y'_j)$. Posons
$\gamma=\alpha\times_S\beta$~; on a par définition
$p\circ\gamma=\alpha\circ p'$, $q\circ\gamma=\beta\circ q'$~; comme
$X'_i=\alpha^{-1}(U_i)$, $Y'_j=\beta^{-1}(V_j)$, on a aussi
${p'}^{-1}(X'_i)=\gamma^{-1}(p^{-1}(U_i))$,
${q'}^{-1}(Y'_j)=\gamma^{-1}(q^{-1}(V_j))$, d'où
\[
  {p'}^{-1}(X'_i)\cap{q'}^{-1}(Y'_j)
  =\gamma^{-1}\bigl(p^{-1}(U_i)\cap q^{-1}(V_j)\bigr)
  =\gamma^{-1}(U_i\times_S V_j)
\]
on conclut comme dans la première partie du raisonnement.

Supposons donc $X$, $Y$, $S$ affines, et soient $B$, $C$, $A$ leurs anneaux
respectifs. Alors $B$ et $C$ sont des $A$-algèbres, $X'$ et $Y'$ des schémas
affines dont les anneaux sont des algèbres quotients $B'$, $C'$ de $B$ et
$C$ respectivement. En outre, on a
$\alpha=({}^a\rho,\widetilde{\rho})$, $\beta=({}^a\sigma,\widetilde{\sigma})$,
où $\rho$ et $\sigma$ sont respectivement les homomorphismes canoniques
$B\to B'$, $C\to C'$ (\hyperref[I.1.7.3-fr]{1.7.3}). Cela étant, on sait que
$X\times_S Y$ (resp. $X'\times_S Y'$) est un schéma affine d'anneau
$B\otimes_A C$ (resp. $B'\otimes_A C'$), et
$\alpha\times_S\beta=({}^a\tau,\widetilde{\tau})$, où $\tau$ est
l'homomorphisme $\rho\otimes\sigma$ de $B\otimes_A C$ dans
$B'\otimes_A C'$ (\hyperref[I.3.2.2-fr]{3.2.2} et
\hyperref[I.3.2.3-fr]{3.2.3})~; comme cet homomorphisme est surjectif,
$\alpha\times_S\beta$ est bien une immersion. En outre, si $\mathfrak{b}$
(resp. $\mathfrak{c}$) est le noyau de $\rho$ (resp. $\sigma$), le noyau de
$\tau$ est $u(\mathfrak{b})+v(\mathfrak{c})$, où $u$ (resp. $v$) est
l'homomorphisme $b\mapsto b\otimes1$ (resp. $c\mapsto1\otimes c$). Comme
$p=({}^a u,\widetilde{u})$ et $q=({}^a v,\widetilde{v})$, ce noyau correspond,
dans le spectre premier de $B\otimes_A C$, à l'ensemble fermé
$p^{-1}(X')\cap q^{-1}(Y')$ (\hyperref[I.1.2.2.1-fr]{1.2.2.1} et
\hyperref[I.1.1.2-fr]{1.1.2 (iii)}), ce qui achève la démonstration.

\begin{corollary}[4.3.2]
\label{I.4.3.2-fr}
Si $f:X\to Y$ est une immersion (resp. une immersion ouverte, une immersion
fermée) et un $S$-morphisme, $f_{(S')}$ est une immersion (resp. une immersion
ouverte, une immersion fermée) pour toute extension $S'\to S$ du préschéma de
base.
\end{corollary}

\subsection{Image réciproque d'un sous-préschéma.}
\label{subsection:I.4.4-fr}

\begin{proposition}[4.4.1]
\label{I.4.4.1-fr}
Soient $f:X\to Y$ un morphisme, $Y'$ un sous-préschéma (resp. un
sous-préschéma fermé, un préschéma induit sur un ouvert) de $Y$, $j:Y'\to Y$
le morphisme d'injection. Alors la projection $p:X\times_Y Y'\to X$ est une
immersion (resp. une immersion fermée, une immersion ouverte)~; le
sous-préschéma de $X$ associé à $p$ a $f^{-1}(Y')$ pour espace sous-jacent~;
en outre, si $j'$ est le morphisme d'injection de ce sous-préschéma dans $X$,
pour qu'un morphisme $h:Z\to X$ soit tel que $f\circ h:Z\to Y$ soit majoré
par $j$, il faut et il suffit que $h$ soit majoré par $j'$.
\end{proposition}

Comme $p=1_X\times_Y j$ (\hyperref[I.3.3.4-fr]{3.3.4}), la première assertion
résulte de (\hyperref[I.4.3.1-fr]{4.3.1})~; la seconde est un cas particulier
de (\hyperref[I.3.5.10-fr]{3.5.10}) (où il faut échanger les rôles de $X$ et
de $Y'$). Enfin, si on a $f\circ h=j\circ h'$, où $h'$ est un morphisme
$Z\to Y'$, il résulte de la définition du produit que l'on a $h=p\circ u$,
où $u$ est un morphisme $Z\to X\times_Y Y'$, d'où la dernière assertion.

Nous dirons que le sous-préschéma de $X$ ainsi défini est l'\emph{image
réciproque} du sous-préschéma $Y'$ de $Y$ par le morphisme $f$, terminologie
qui s'accorde avec celle introduite
\oldpage[I]{126}
de façon plus générale dans (\hyperref[I.3.3.6-fr]{3.3.6}). Quand nous
parlerons de $f^{-1}(Y')$ comme d'un sous-préschéma de $X$, c'est toujours de
ce sous-préschéma qu'il s'agira.

Lorsque les préschémas $f^{-1}(Y')$ et $X$ sont identiques, $j'$ est
l'identité et tout morphisme $h:Z\to X$ est donc majoré par $j'$, donc
\emph{le morphisme $f:X\to Y$ se factorise alors en}
$X\xrightarrow{g}Y'\xrightarrow{j}Y$.

Lorsque $y$ est un point \emph{fermé} de $Y$ et
$Y'=\operatorname{Spec}(k(y))$ le plus petit sous-préschéma fermé de $Y$ ayant
$\{y\}$ comme espace sous-jacent (\hyperref[I.4.1.9-fr]{4.1.9}), le
sous-préschéma fermé $f^{-1}(Y')$ est canoniquement isomorphe à la
\emph{fibre} $f^{-1}(y)$ définie en (\hyperref[I.3.6.2-fr]{3.6.2}), à
laquelle on l'identifiera.

\begin{corollary}[4.4.2]
\label{I.4.4.2-fr}
Soient $f:X\to Y$, $g:Y\to Z$ deux morphismes, $h=g\circ f$ leur composé.
Pour tout sous-préschéma $Z'$ de $Z$, les sous-préschémas
$f^{-1}(g^{-1}(Z'))$ et $h^{-1}(Z')$ de $X$ sont identiques.
\end{corollary}

Cela résulte de l'existence de l'isomorphisme canonique
$X\times_Y(Y\times_Z Z')\xrightarrow{\sim}X\times_Z Z'$
(\hyperref[I.3.3.9.1-fr]{3.3.9.1}).

\begin{corollary}[4.4.3]
\label{I.4.4.3-fr}
Soient $X'$, $X''$ deux sous-préschémas de $X$, $j':X'\to X$,
$j'':X''\to X$ les morphismes d'injection~; alors, ${j'}^{-1}(X'')$ et
${j''}^{-1}(X')$ sont tous deux égaux à la borne inférieure
$\inf(X',X'')$ de $X'$ et $X''$ pour la relation d'ordre entre
sous-préschémas, et canoniquement isomorphée à $X'\times_X X''$.
\end{corollary}

Cela résulte aussitôt de (\hyperref[I.4.4.1-fr]{4.4.1}) et
(\hyperref[I.4.1.10-fr]{4.1.10}).

\begin{corollary}[4.4.4]
\label{I.4.4.4-fr}
Soient $f:X\to Y$ un morphisme, $Y'$, $Y''$ deux sous-préschémas de $Y$~;
on a
\[
  f^{-1}(\inf(Y',Y''))=\inf(f^{-1}(Y'),f^{-1}(Y'')).
\]
\end{corollary}

Cela résulte de l'existence de l'isomorphisme canonique entre
$(X\times_Y Y')\times_X(X\times_Y Y'')$ et
$X\times_Y(Y'\times_Y Y'')$ (\hyperref[I.3.3.9.1-fr]{3.3.9.1}).

\begin{proposition}[4.4.5]
\label{I.4.4.5-fr}
Soient $f:X\to Y$ un morphisme, $Y'$ un sous-préschéma fermé de $Y$ défini
par un faisceau quasi-cohérent $\mathscr{K}$ d'idéaux de $\mathscr{O}_Y$
(\hyperref[I.4.1.3-fr]{4.1.3})~; le sous-préschéma fermé $f^{-1}(Y')$ de $X$
est alors défini par le faisceau quasi cohérent d'idéaux
$f^*(\mathscr{K})\mathscr{O}_X$ de $\mathscr{O}_X$.
\end{proposition}

La question est évidemment locale sur $X$ et $Y$~; il suffit alors de
remarquer que si $B$ est une $A$-algèbre et $\mathfrak{K}$ un idéal de $B$,
on a $A\otimes_B(B/\mathfrak{K})=A/\mathfrak{K}A$, et d'appliquer
(\hyperref[I.1.6.9-fr]{1.6.9}).

\begin{corollary}[4.4.6]
\label{I.4.4.6-fr}
Soient $X'$ un sous-préschéma fermé de $X$ défini par un faisceau quasi
cohérent d'idéaux $\mathscr{J}$ de $\mathscr{O}_X$, $i$ l'injection
$X'\to X$~; pour que la restriction $f\circ i$ de $f$ à $X'$ soit majorée
par l'injection $j:Y'\to Y$ (autrement dit se factorise en $j\circ g$, où
$g$ est un morphisme $X'\to Y'$), il faut et il suffit que
$f^*(\mathscr{K})\mathscr{O}_X\subset\mathscr{J}$.
\end{corollary}

Il suffit d'appliquer à $i$ la prop. (\hyperref[I.4.4.1-fr]{4.4.1}) en
tenant compte de (\hyperref[I.4.4.5-fr]{4.4.5}).

\subsection{Immersions locales et isomorphismes locaux.}
\label{subsection:I.4.5-fr}

\begin{definition}[4.5.1]
\label{I.4.5.1-fr}
Soit $f:X\to Y$ un morphisme de préschémas. On dit que $f$ est une
\emph{immersion locale en un point $x\in X$} s'il existe un voisinage ouvert
$U$ de $x$ dans $X$ et un voisinage ouvert $V$ de $f(x)$ dans $Y$ tels que
la restriction de $f$ au préschéma induit $U$ soit une immersion fermée de
$U$ dans le préschéma induit $V$. On dit que $f$ est une immersion locale si
$f$ est une immersion locale en tout point de $X$.
\end{definition}

\begin{definition}[4.5.2]
\label{I.4.5.2-fr}
On dit qu'un morphisme $f:X\to Y$ est un isomorphisme local en
\oldpage[I]{127}
un point $x\in X$ s'il existe un voisinage ouvert $U$ de $x$ dans $X$ tel
que la restriction de $f$ au préschéma induit $U$ soit une immersion ouverte
de $U$ dans $Y$. On dit que $f$ est un isomorphisme local si $f$ est un
isomorphisme local en tout point de $X$.
\end{definition}

\begin{env}[4.5.3]
\label{I.4.5.3-fr}
Une immersion (resp. une immersion fermée) $f:X\to Y$ peut donc se
caractériser comme une immersion locale telle que $f$ soit un homéomorphisme
de l'espace sous-jacent à $X$ sur une partie (resp. une partie fermée) de
$Y$. Une immersion ouverte $f$ peut se caractériser comme un isomorphisme
local \emph{injectif}.
\end{env}

\begin{proposition}[4.5.4]
\label{I.4.5.4-fr}
Soient $X$ un préschéma irréductible, $f:X\to Y$ un morphisme injectif
dominant. Si $f$ est une immersion locale, $f$ est une immersion et $f(X)$
est ouvert dans $Y$.
\end{proposition}

En effet, soit $x\in X$ et soient $U$ un voisinage ouvert de $x$, $V$ un
voisinage ouvert de $f(x)$ dans $Y$ tels que la restriction de $f$ à $U$
soit une immersion fermée dans $V$~; comme $U$ est dense dans $X$, $f(U)$
est dense dans $Y$ par hypothèse, donc $f(U)=V$ et $f$ est un homéomorphisme
de $U$ sur $V$~; l'hypothèse que $f$ est injectif entraîne que
$f^{-1}(V)=U$, d'où aussitôt la proposition.

\begin{proposition}[4.5.5]
\label{I.4.5.5-fr}
\begin{enumerate}
  \item[\rm(i)] Le composé de deux immersions locales (resp. de deux
    isomorphismes locaux) est une immersion locale (resp. un isomorphisme
    local).
  \item[\rm(ii)] Soient $f:X\to X'$, $g:Y\to Y'$ deux $S$-morphismes. Si
    $f$ et $g$ sont des immersions locales (resp. des isomorphismes locaux),
    il en est de même de $f\times_S g$.
  \item[\rm(iii)] Si un $S$-morphisme $f$ est une immersion locale (resp. un
    isomorphisme local), il en est de même de $f_{(S')}$ pour toute extension
    $S'\to S$ du préschéma de base.
\end{enumerate}
\end{proposition}

D'après (\hyperref[I.3.5.1-fr]{3.5.1}), il suffit de prouver (i) et (ii).

(i) résulte aussitôt de la transitivité des immersions fermées (resp.
ouvertes) (\hyperref[I.4.2.4-fr]{4.2.4}) et du fait que si $f$ est un
homéomorphisme de $X$ sur une partie fermée de $Y$, pour tout ouvert
$U\subset X$, $f(U)$ est ouvert dans $f(X)$, donc il existe une partie
ouverte $V$ de $Y$ telle que $f(U)=V\cap f(X)$, et $f(U)$ est par suite fermé
dans $V$.

Pour démontrer (ii), soient $p$, $q$ les projections de $X\times_S Y$, $p'$,
$q'$ celles de $X'\times_S Y'$. Il existe par hypothèse des voisinages
ouverts $U$, $U'$, $V$, $V'$ de $x=p(z)$, $x'=p'(z')$, $y=q(z)$,
$y'=q'(z')$ respectivement, tels que les restrictions de $f$ et $g$ à $U$
et $V$ respectivement soient des immersions fermées (resp. ouvertes) dans
$U'$ et $V'$ respectivement. Comme l'espace sous-jacent de $U\times_S V$ et
celui de $U'\times_S V'$ s'identifient aux voisinages ouverts
$p^{-1}(U)\cap q^{-1}(V)$ et ${p'}^{-1}(U')\cap{q'}^{-1}(V')$ de $z$ et
$z'$ respectivement (\hyperref[I.3.2.7-fr]{3.2.7}), la proposition résulte
de (\hyperref[I.4.3.1-fr]{4.3.1}).
